{\footnotesize
\setlength{\tabcolsep}{3pt}
\renewcommand{\arraystretch}{1.12}
\begin{longtable}{>{\raggedright\arraybackslash}p{2.6cm}>{\raggedright\arraybackslash}p{3.7cm}>{\raggedright\arraybackslash}p{5.7cm}rrrrr}
\caption{Results across the cases of Contraste: for each case the question it asks, the headline result read from its committed artifacts with the uncertainty the artifact carries, the opinion level of its validation reports (Op., from 1 to 4 on the scale of the euro-area supervisor) and the findings it raised by severity S1 (critical) to S4 (minor). Case C19 has no committed result.}\label{tab:results-summary}\\
\toprule
Case & Question & Headline result and its uncertainty & Op. & S1 & S2 & S3 & S4 \\
\midrule
\endfirsthead
\toprule
Case & Question & Headline result and its uncertainty & Op. & S1 & S2 & S3 & S4 \\
\midrule
\endhead
\bottomrule
\endfoot
\textbf{C01} Retail cards PD: champion vs challenger & Does a monotone WoE scorecard match the challengers the literature says can beat it, once calibration is required, and what do the validation tests say when the population shifts? & Scorecard AUC \csHoScAuc{} (95\,\% CI \csHoScAucLo{} to \csHoScAucHi{}) on \csHoN{} held-out cards; LightGBM adds \csHoLgDlDiff{} (DeLong $p$ \csHoLgDlP{}). & 3 & 0 & 31 & 45 & 7 \\
\textbf{C02} Corporate PD across 1 to 5 year horizons & How fast does discrimination fall from the 1-year to the 5-year horizon, do the learned rungs beat the scorecard by a margin a paired test supports, and do they stay calibrated? & One-year file: scorecard AUC \csCtwoAScAuc{} (95\,\% CI \csCtwoAScAucLo{} to \csCtwoAScAucHi{}); the explainable boosting machine adds \csCtwoAEbDiff{} (\csCtwoAEbDiffLo{} to \csCtwoAEbDiffHi{}). & 3 & 0 & 2 & 20 & 18 \\
\textbf{C03} SME PD by vintage through 2008 & Does a PD model built on approval cohorts old enough to have been observed for five years stay calibrated for the cohorts that lived through 2008-2010, and what does a validator find in the cohorts before and after? & FY2008--10 cohort: scorecard AUC 0.815 (95\,\% CI 0.812 to 0.818), LightGBM 0.832 (0.830 to 0.835); observed 5-year PD 13.1\,\% against 6.5\,\% predicted. & 3 & 0 & 32 & 28 & 5 \\
\textbf{C04} Rating transitions and TTC PD by grade & What do the agencies' published migrations say about PDs by grade, and how far can a migration matrix be trusted to give them, project them, or show the chain is Markov and stable? & Long-run PD of CCC--C \rtgLraSpCCCC{}\,\% (S\&P), \rtgLraMoodysCCCC{}\,\% (Moody's), \rtgLraFitchCCCC{}\,\% (Fitch); AAA of S\&P: \rtgLradSpAAA{} defaults in \rtgLranSpAAA{} rating-years, 95\,\% Jeffreys upper bound \rtgJeffupSpAAA{}\,\%. & 3 & 0 & 17 & 27 & 3 \\
\textbf{C05} Low-default portfolios and PD calibration & How should a PD curve be carried to a new year, and what can be said about PDs with almost no defaults? & IRB risk weight of one book under 9 PD-curve approaches for 2010: 55.0\,\% to 98.5\,\% of EAD (exact recomputation, no sampling interval). & 3 & 0 & 7 & 3 & 5 \\
\textbf{C06} Lifetime PD, IFRS 9 staging and ECL on public small-business loans & How do lifetime PD, IFRS 9 staging and expected credit loss behave on public small-business loans across reporting dates, and what does a validator find when the model of each date is tested on what happened after it? & Cox c-index \pvCindexCoxMin{} to \pvCindexCoxMax{} and PC-Hazard c-index \pvCindexPchMin{} to \pvCindexPchMax{} over \pvDatesTotal{} reporting dates (range over dates, no interval). & 3 & 0 & 14 & 6 & 32 \\
\textbf{C07} Chile provisions trilogy: CMF B-1 vs IFRS 9 vs BdE Annex 9 & How far apart are the provisions of CMF B-1, IFRS 9 and BdE Annex 9 on the same portfolio under stress? & Base scenario provision (thousand million CLP): CMF B-1 152.0, IFRS 9 97.6, Banco de Espa\~na Annex 9 79.6; IFRS 9 over B-1 0.64 (synthetic bank, no interval). & 3 & 0 & 3 & 14 & 12 \\
\textbf{C08} LGD and downturn LGD & How does loss given default move across default cohorts, and what does a downturn LGD add? & Calibration by segment red in \LE{cal-red} of \LE{cal-total} cohort-by-rung tests, \LE{cal-positive} gaps under-predicting; $R^2$ between \LE{r2-min} and \LE{r2-max}; gradient boosting has the lowest RMSE in \LE{bestcount-L5-gradient-boosted} of 6 cohorts. & 3 & 0 & 17 & 6 & 0 \\
\textbf{C09} EAD and CCF for revolving lines & Can exposure at default and the credit conversion factor of revolving lines be recovered when utilisation accelerates? & Untreated pooled CCF bias against the known 0.714: fixed horizon $-$0.019 [$-$0.026, $-$0.012], cohort $-$0.223 [$-$0.231, $-$0.215], variable horizon $-$0.150 [$-$0.158, $-$0.142] (95\,\% intervals). & 3 & 0 & 13 & 21 & 9 \\
\textbf{C10} Portfolio credit capital: ASRF, Vasicek fit, copula MC & Do the ASRF formula, a fitted Vasicek model and a copula simulation agree on portfolio credit capital, including the 2007 to 2012 tail? & Fitted asset correlation \ccRhoMlRes{} for residential real estate (\ccRhoMlLoRes{} to \ccRhoMlHiRes{}) and \ccRhoMlCre{} for commercial real estate (\ccRhoMlLoCre{} to \ccRhoMlHiCre{}), 95\,\% profile intervals. & 3 & 0 & 9 & 13 & 14 \\
\textbf{C11} Market VaR/ES backtesting across 2008, 2020, 2022, 2023 & Do VaR and expected shortfall models pass their backtests across the 2008, 2020, 2022 and 2023 episodes? & Eligible models of \mkNModels{}: \mkEligCalm{} in the calm window, \mkEligGfc{} in 2007--09, \mkEligCovid{} in 2020, \mkEligRates{} in the 2022 rates window, \mkEligFx{} in the FX window (Kupiec, conditional coverage and Acerbi--Szekely tests). & 3 & 0 & 31 & 2 & 1 \\
\textbf{C12} How large is the counterparty exposure of a swap book, and what CVA does it imply? & How large is the counterparty exposure of a swap book, what CVA does it imply, and would the simulated exposures have held against the curves later realised? & Netting-set EEPE USD \ccrEepeNetting{}; one standard error on the mean reversion (\ccrHwA{} $\pm$ \ccrHwASe{}) moves it by \ccrEepeAPlusPct{}\,\% to +\ccrEepeAMinusPct{}\,\%. & 3 & 0 & 27 & 0 & 0 \\
\textbf{C13} FRTB desk eligibility: PLA and desk backtesting & Do the P\&L attribution and desk backtesting tests of FRTB detect planted model deficiencies? & Profit-and-loss attribution flags the missing-basis model in 96.0\,\% (94.6 to 97.0) of 1{,}000 replications and the stale-volatility model in 0.0\,\% (0.0 to 0.4). & 3 & 0 & 4 & 8 & 7 \\
\textbf{C14} IRRBB, six shocks on a USD balance sheet: the bank that failed in 2023 & Would the standardised EVE outlier test and the NII sensitivity have flagged the position of the bank that failed in March 2023, read on its call reports, and how much does the answer hang on the deposit assumptions? & EVE risk measure at 2022-Q4 USD 3.07 thousand million with behavioural deposits, 10.98 contractual; 18.1\,\% and 64.6\,\% of Tier 1 capital (no interval). & 4 & 3 & 4 & 5 & 0 \\
\textbf{C15} IRRBB on a CLP and UF balance sheet: RAN 21-13 Annex 1 & How much EVE and NII does a synthetic Chilean bank lose under the CMF's prescribed shocks on the Banco Central curves, how much of it is the indexation of its UF position, and what does the 50 percent offset between the pesos and UF legs change? & Long-UF variant: $\Delta$EVE 93{,}816 and $\Delta$NII 11{,}213 million CLP; the indexation charge is \pxSevenShare{}\,\% of the $\Delta$NII of the short-UF variant (no interval). & 3 & 0 & 1 & 17 & 6 \\
\textbf{C16} Liquidity stress, LCR and the March 2023 run & Would the LCR and liquidity stress models have flagged the March 2023 deposit run? & Survival period 62 days under the adverse run-off; the run-off multiplier that brings it to \pxElevenTarget{} days is \pxElevenAdverse{} (no interval). & 3 & 0 & 13 & 5 & 1 \\
\textbf{C17} Operational risk: LDA vs SMA & How do the loss distribution approach and the standardised measurement approach differ on heavy-tailed operational losses? & Tail index 0.8, 10 years: LDA capital EUR 148.5 million against SMA 142.0; the interquartile range of the LDA capital is \pxTwelveIqrMin{} to \pxTwelveIqrMax{} times the truth over \pxTwelveHistories{} histories. & 3 & 0 & 13 & 8 & 0 \\
\textbf{C18} ICAAP and IAPE capital aggregation under macro scenarios & Does Banco Andino hold its internal capital target for three years under each scenario once credit, market, IRRBB and operational risk are aggregated, and how much of the answer rests on the aggregation method? & Severely adverse scenario: CET1 ratio 7.36\,\% of RWA at the end of year 3; smallest headroom over the internal target $-$1.25 points; the band of aggregation methods is shown with the headline capital (no interval). & 4 & 2 & 6 & 17 & 8 \\
\textbf{C19} Supervisory stress model replication (known answer) & Supervisory stress replication against a published answer. & No committed result: the case is planned and the report states none. & -- & -- & -- & -- & -- \\
\textbf{C20} ML model validation: drift, robustness, explanation stability & Do performance, drift, attribution and reason-code stability survive shocks the model has and has not seen, shifted populations, input noise and retraining seeds, and does the leakage screen catch a planted leak? & Largest decision-flip share over the applied shocks 1.1\,\% on the stressed German credit sample; no sampling interval is committed. & 3 & 0 & 18 & 9 & 12 \\
\textbf{C21} Fair-lending disparity testing & Do disparity tests find the same gaps in lending decisions from one year to the next? & 2023--24 pair: approval rate of Black or African American applicants 63.2\,\% against 78.0\,\% for White applicants, a difference of $-$14.8 points (95\,\% Wald interval $-$16.2 to $-$13.4). & 3 & 0 & 6 & 14 & 18 \\
\textbf{C22} Validating the validator: size and power of every test & How often does each validation test reject a right model, and which defects does it catch? & 19 of \pxOneTests{} tests hold their size at 5 and 1\,\%; Hosmer--Lemeshow with $G-2$ degrees of freedom rejects 0.110 of null validation samples (Monte Carlo bound 0.061); \pxOneRates{} rates carry their standard error. & 3 & 0 & 18 & 9 & 6 \\
\end{longtable}
}
